\Cref{tab:transfer} reports the results on LSFB-ISOL with all training data. Trained from scratch,
the Conv-Transformer reaches \numLsfbScratchTest{}\,\% top-1 and \numLsfbScratchTestFOne{}\,\%
macro-F1 on the 40 official test signers, well above the statistical baselines (\numLsfbLogregTest{}\,\% and
\numLsfbLogregTestFOne{}\,\%). Macro-F1 is markedly lower than top-1 on this corpus because of its long tail: the ten
most frequent glosses account for \numLsfbTopTenShare{}\,\% of the test clips.

\paragraph{Validation and test accuracy.}
On LSFB, validation accuracy exceeds test accuracy by far more than on ASL (\numGapValClip{} against
\numGapTestClip{}\,\%). The validation set is small and unbalanced across signers: three of its
\numGapValSigners{} signers contribute \numGapTopThreeShare{}\,\% of its clips. Averaged per signer,
accuracy is \numGapValPerSigner{}\,$\pm$\,\numGapValPerSignerSem{}\,\% on validation (seed 0) and
\numGapTestPerSigner{}\,$\pm$\,\numGapTestPerSignerSem{}\,\% on test (mean over seeds; mean $\pm$
standard error over signers), and re-weighting the test clips to the class distribution of the validation set raises
test accuracy to \numGapTestReweighted{}\,\%. LSFB clips come from dialogues: test signers whose
conversation partner belongs to the training signers (\numGapTestPartnerInN{} of
\numGapTestSigners{}) are recognised with \numGapTestPartnerIn{}\,\% accuracy, the others with
\numGapTestPartnerOut{}\,\%, which suggests that shared conversations, and therefore shared topics
and signing styles, also favour held-out signers. Validation scores are therefore used only to
compare configurations, never as an estimate of performance on new signers.

\paragraph{ASL pre-training.}
Initialising the encoder from the ASL model does not improve recognition. With all training data,
fine-tuning the ASL model yields \numLsfbFtTest{}\,\% (\numLsfbFtVal{}\,\% on validation, against
\numLsfbScratchVal{}\,\% from scratch), and freezing the encoder during the first five epochs only
narrows the gap (\numLsfbFreezeTest{}\,\%). Pre-training is not merely useless but slightly harmful: relative to training from scratch,
test accuracy changes by \numLsfbFtDiff{} points with fine-tuning (signer-level 95\,\% interval
[\numLsfbFtCiLo{}, \numLsfbFtCiHi{}], \numLsfbFtP{}) and by \numLsfbFreezeDiff{} points with the
frozen variant ([\numLsfbFreezeCiLo{}, \numLsfbFreezeCiHi{}]). The learning curves of \cref{fig:transfer} show that
this result does not depend on the size of the training set. Both initialisations are
indistinguishable from 5 to 25 examples per sign (the largest difference, \numGainFew{} points at
5 examples, lies within one seed standard deviation), and the model trained from scratch leads
from 50 examples per sign onwards, by \numScratchLeadAll{} points with all data.

\paragraph{Controls.}
Two controls exclude an inadequate fine-tuning procedure as the explanation. First, a linear probe
trained on top of the frozen ASL encoder with all LSFB training data reaches \numProbeAll{}\,\%
validation top-1, far below the model trained from scratch (\numLsfbScratchVal{}\,\%) and even
below the logistic regression on temporal statistics (\numLsfbLogregVal{}\,\%). The ASL
representation therefore retains some information about LSFB signs, but less than simple features
of the target data. Second, with 25 examples per sign, fine-tuning with a five times lower
learning rate gives \numLowlrTwentyFive{}\,\% validation top-1, below both standard fine-tuning
(\numFtTwentyFive{}\,\%) and training from scratch (\numScratchTwentyFive{}\,\%) with the same seed.
Both controls use a single seed, and every transfer run starts from the ASL model of seed 0, so
the variability of the pre-trained model itself is not sampled. During the linear probe, the
batch-normalisation statistics of the frozen encoder were still updated on LSFB, a mild form of
adaptation; the probe is therefore, if anything, favoured, which does not change its conclusion.

\begin{figure}[t]
  \centering
  \begin{minipage}[b]{0.44\linewidth}
    \includegraphics[width=\linewidth]{transfer_curve}
    \caption{LSFB validation top-1 as a function of the number of training examples per sign
    (mean $\pm$ s.d.\ over three seeds). ``All'' corresponds to a median of \numLsfbMedianPerClass{} examples per
    sign.}
    \label{fig:transfer}
  \end{minipage}\hfill
  \begin{minipage}[b]{0.53\linewidth}
    \centering
    \small
    \captionof{table}{LSFB-ISOL, 100 glosses, all training data. Mean $\pm$ s.d.\ over three
    seeds, values in \%.}
    \label{tab:transfer}
    \resizebox{\linewidth}{!}{\begin{tabular}{lcccc}
\toprule
 & Validation & \multicolumn{3}{c}{Test (40 unseen signers)} \\
\cmidrule(lr){2-2}\cmidrule(lr){3-5}
Model & Top-1 & Top-1 & Top-5 & Macro-F1 \\
\midrule
Chance & 1.0 & 1.0 & 5.0 & -- \\
Statistics + log.\ regression & 57.0 & 49.0 & 75.3 & 40.9 \\
Statistics + LightGBM & 54.8 & 47.1 & 72.6 & 37.3 \\
\midrule
Conv-Transformer, from scratch & 77.0 $\pm$ 0.3 & 68.3 $\pm$ 0.2 & 87.7 $\pm$ 0.1 & 63.9 $\pm$ 0.2 \\
\quad ASL-pretrained, fine-tuned & 74.9 $\pm$ 0.4 & 66.2 $\pm$ 0.2 & 87.1 $\pm$ 0.4 & 61.7 $\pm$ 0.3 \\
\quad ASL-pretrained, encoder frozen 5 epochs & 75.6 $\pm$ 0.6 & 66.6 $\pm$ 0.2 & 86.5 $\pm$ 0.5 & 62.0 $\pm$ 0.3 \\
\bottomrule
\end{tabular}
}
  \end{minipage}
\end{figure}
